\section{The global distribution of zero-velocity lattice parameters}
\label{sec:velocity-zero-law}

The same singularity geometry has a second interpretation: one may fix
the index and vary the lattice spacing.  Define
\[
 p_n(L)=\sum_{j=1}^n\stirling{n}{n-j+1}\frac{(-L)^j}{j!},
 \qquad v_n(A)=A^{-n}p_n(\log A).
\]
For $n\geq2$, write the positive zeros of $p_n$ in increasing order as
$L_{n,1}<\cdots<L_{n,n-1}$, and put $A_{n,j}=e^{L_{n,j}}$.
Thus $A_{n,j}-1$ is a lattice spacing where the initial diagonal
velocity vanishes.  Equivalently, these $L_{n,j}$ are the positive
zeros of Cohen's Lambert polynomial $L_n(L)=(-1)^np_n(L)$
\cite[Proposition~4.8]{Cohen2021}; the law below is therefore a bulk
zero law for that known family.

\begin{lemma}[Real-rootedness and compact uniformity]
\label{lem:zero-law-inputs}
The polynomial $p_n$ has exactly the simple zero $0$ and $n-1$ simple
positive zeros.  The additive asymptotic
\eqref{eq:velocity-leading-cosine} is uniform for $A$ in any compact
subinterval of $(1,\infty)$.
\end{lemma}

\begin{proof}
The real-rootedness is the elementary diagonal result from
\cite{HarmonicBaseline}; here is its operator proof.  With $D=d/dx$,
\[
 n p_n(L)=
 \left.\left(\prod_{j=1}^{n-1}(1+D/j)\right)x^n\right|_{x=-L}.
\]
For a polynomial $P$ with real roots $r_i$ of multiplicities $m_i$,
the roots of $P+cP'$ away from inherited roots solve
$\sum_i m_i/(x-r_i)=-1/c$.  For $c>0$, the left side is strictly
decreasing between its poles.  There is one new simple root in each
interval between distinct roots and one to the left of the smallest
root; an inherited root loses one multiplicity.  Starting with $x^n$
and applying $n-1$ such operators gives $n$ simple nonpositive roots,
of which zero is the largest.  Replacing $x$ by $-L$ proves the assertion.

For uniformity, fix a compact parameter interval.  Its angles stay a
positive distance from $0$ and $\pi$, its values of $a$ stay positive,
and the two critical points remain separated.  The critical parameters
and local coefficients depend real analytically on $A$.  In the
normalized variable $r=t/R(A)$, the rectangle estimates used in
Theorem~\ref{thm:selected-sheet} have uniform strict margins away from
small neighborhoods of the two equality points.  The right boundary
can be chosen with one common sufficiently large $M$.  On the compact
set of regular boundary parameter pairs, the analytic implicit function
theorem therefore supplies a uniform collar after taking a finite cover.
At the two moving endpoints the normalized local equation
\eqref{eq:local-implicit-equation} has a quadratic coefficient bounded
away from zero, giving uniform Puiseux neighborhoods and remainder
bounds, again by compactness.  The resulting transfer contours may be
chosen with fixed normalized sizes and fixed angular dents.  The proof
of Theorem~\ref{thm:velocity-asymptotics} then gives a remainder constant
independent of $A$ in the chosen interval.
\end{proof}

\begin{theorem}[Limiting zero law and explicit bulk quantiles]
\label{thm:velocity-zero-law}
For every $A>1$,
\begin{equation}\label{eq:zero-law-cdf}
 \lim_{n\to\infty}\frac1{n-1}
       \#\{j:A_{n,j}\leq A\}=\frac{\theta(A)}\pi.
\end{equation}
Consequently, if $j_n/(n-1)\to x\in(0,1)$, then
\begin{equation}\label{eq:zero-law-quantile}
 A_{n,j_n}\longrightarrow
 \frac{\pi x}{\sin(\pi x)}\exp\bigl(1-\pi x\cot(\pi x)\bigr).
\end{equation}
In the logarithmic coordinate $L=\log A$, the limiting probability
measure has density
\begin{equation}\label{eq:zero-law-density}
 \frac{d\mu}{dL}
   =\frac{\theta}{\pi(a^2+\theta^2)},\qquad
 L=1-\theta\cot\theta+\log\frac\theta{\sin\theta},
 \quad a=1-\theta\cot\theta.
\end{equation}
Equivalently, in the original coordinate $A>1$ its density is
\begin{equation}\label{eq:zero-law-density-A}
 \frac{d\mu}{dA}=\frac{\theta(A)}
              {\pi A\bigl(a(A)^2+\theta(A)^2\bigr)}.
\end{equation}
\end{theorem}

\begin{proof}
Let $0<\alpha<\beta<\pi$.  The parameter $A$ is a strictly increasing
function of $\theta$.  Uniformly on $[\alpha,\beta]$, the polynomial
$p_n(L(\theta))$ has the sign of
\[
 \cos(n\theta+\phi(\theta))+O_{\alpha,\beta}(1/n),
\]
apart from the ambiguity allowed by a small cosine.  The smooth function
$\phi$ has bounded derivative there, so $n\theta+\phi(\theta)$ is
strictly increasing for all sufficiently large $n$.  At its successive
integer multiples of $\pi$, the leading cosines are alternately $1$
and $-1$, and the error is eventually smaller than $1/2$.  Between
successive such points the intermediate value theorem supplies a
polynomial zero.  The number of these successive pairs is
$n(\beta-\alpha)/\pi+O_{\alpha,\beta}(1)$.  Hence the lower limit
of the fraction of zeros whose angle lies in $[\alpha,\beta]$ is at
least $(\beta-\alpha)/\pi$.

Fix $\theta_0\in(0,\pi)$.  Applying this lower bound on
$[\epsilon,\theta_0]$ gives a lower limit at least
$(\theta_0-\epsilon)/\pi$ for the fraction below $\theta_0$.
Applying it on $[\theta_0+\epsilon,\pi-\epsilon]$ and subtracting
from the total of exactly $n-1$ positive roots gives an upper limit
at most $(\theta_0+2\epsilon)/\pi$.  Letting $\epsilon$ decrease
to zero proves \eqref{eq:zero-law-cdf}.  Strict monotonicity and
continuity of its limiting distribution function give the quantile
limit \eqref{eq:zero-law-quantile} by bracketing $x$ between nearby
continuity points.

Finally,
\[
 \frac{dL}{d\theta}
   =\theta\csc^2\theta+\frac1\theta-2\cot\theta
   =\frac{(1-\theta\cot\theta)^2+\theta^2}{\theta}
   =\frac{a^2+\theta^2}{\theta}.
\]
Pushing forward the uniform measure $d\theta/\pi$ proves
\eqref{eq:zero-law-density}.  Since $\theta$ runs from $0$ to $\pi$,
the density has total mass one, and the CDF also proves tightness of the
empirical measures.
\end{proof}

This distribution has a heavy logarithmic tail.  As $L\to\infty$,
$a\sim L$ and $\theta\to\pi$, so $d\mu/dL\sim L^{-2}$ and
$\mu((L,\infty))\sim L^{-1}$.  There is no contradiction with the
exact finite-degree first moment
\begin{equation}\label{eq:zero-law-exact-moment}
 \sum_{j=1}^{n-1}L_{n,j}=nH_{n-1},
\end{equation}
obtained by comparing the leading two coefficients of $p_n$: the
limiting distribution has infinite mean.  The bulk law concerns indices
proportional to $n$; it gives no uniform estimate for either extreme
root, whose scaling requires separate analysis.
